\section{Symulacja i testy}
\label{sec:testy}

\subsection{Metodyka}
Działanie układu sprawdzono w symulatorze LOGO!Soft Comfort. Wejście
$I_1$ skonfigurowano jako przycisk chwilowy, a czas przytrzymania
mierzono stoperem. Przed pierwszym testem układ znajdował się w stanie
$S_0$ (wszystkie przerzutniki z parametrem \texttt{Rem=off}).
Scenariusze wykonywano jeden po drugim, a stan następny jednego testu
był stanem początkowym kolejnego.

\subsection{Scenariusze testowe}
Tabela~\ref{tab:testy} zawiera sekwencję testów obejmującą wszystkie
przejścia z tabeli~\ref{tab:przejscia}. Wynik każdego testu odczytano
ze stanu wyjść $Q_1$ i $Q_2$ oraz markera $M_3$.

\begin{table}[H]
\centering
\caption{Scenariusze testowe i oczekiwane wyniki.}
\label{tab:testy}
\begin{tabular}{rllcll|cc}
\toprule
Nr & Stan przed & Akcja & $t$ [s] & Stan po & & $Q_2$ & $Q_1$ \\
\midrule
 1 & $S_0$    & K & 0,2 & $S_0$    & brak reakcji           & 0 & 0 \\
 2 & $S_0$    & Ś & 1   & $S_0$    & brak reakcji           & 0 & 0 \\
 3 & $S_0$    & D & 3   & $S_2$    & uruchomienie telefonu  & 1 & 1 \\
 4 & $S_2$    & K & 0,2 & $S_1$    & ekran wygaszony        & 1 & 0 \\
 5 & $S_1$    & K & 0,2 & $S_2$    & ekran włączony         & 1 & 1 \\
 6 & $S_2$    & Ś & 1   & $S_{3b}$ & otwarcie menu          & 1 & 1 \\
 7 & $S_{3b}$ & K & 0,2 & $S_2$    & anulowanie menu        & 1 & 1 \\
 8 & $S_2$    & K & 0,2 & $S_1$    & ekran wygaszony        & 1 & 0 \\
 9 & $S_1$    & Ś & 1   & $S_{3a}$ & menu wymusza ekran     & 1 & 1 \\
10 & $S_{3a}$ & K & 0,2 & $S_1$    & ekran znów wygaszony   & 1 & 0 \\
11 & $S_1$    & Ś & 1   & $S_{3a}$ & otwarcie menu          & 1 & 1 \\
12 & $S_{3a}$ & Ś & 1   & $S_{3a}$ & menu pozostaje otwarte & 1 & 1 \\
13 & $S_{3a}$ & D & 3   & $S_0$    & wyłączenie, menu zamknięte & 0 & 0 \\
14 & $S_0$    & D & 3   & $S_2$    & uruchomienie, menu zamknięte & 1 & 1 \\
15 & $S_2$    & Ś & 1   & $S_{3b}$ & otwarcie menu          & 1 & 1 \\
16 & $S_{3b}$ & D & 3   & $S_0$    & wyłączenie telefonu    & 0 & 0 \\
17 & $S_0$    & D & 3   & $S_2$    & uruchomienie telefonu  & 1 & 1 \\
18 & $S_2$    & D & 3   & $S_0$    & wyłączenie telefonu    & 0 & 0 \\
19 & $S_0$    & D & 3   & $S_2$    & uruchomienie telefonu  & 1 & 1 \\
20 & $S_2$    & K & 0,2 & $S_1$    & ekran wygaszony        & 1 & 0 \\
21 & $S_1$    & D & 3   & $S_0$    & wyłączenie telefonu    & 0 & 0 \\
\bottomrule
\end{tabular}
\end{table}

Test~3 sprawdza wymaganie R1 (uruchomienie włącza ekran), test~13
sprawdza zamknięcie menu przy wyłączaniu telefonu, a test~14 -- że po
ponownym uruchomieniu menu nie jest otwarte. Testy 1, 2 potwierdzają,
że przy wyłączonym telefonie kliknięcia K i Ś nie wywołują reakcji.

\subsection{Testy graniczne}
Dodatkowo sprawdzono zachowanie układu w przypadkach granicznych
(tab.~\ref{tab:graniczne}).

\begin{table}[H]
\centering
\caption{Testy graniczne.}
\label{tab:graniczne}
\begin{tabular}{lll}
\toprule
Scenariusz & Oczekiwany wynik & Wymaganie \\
\midrule
trzymanie \SI{10}{\second} w stanie $S_2$ & jedno wyłączenie telefonu, bez migotania $Q_2$ & R8 \\
trzymanie \SI{10}{\second} w stanie $S_0$ & jedno uruchomienie telefonu, bez migotania $Q_2$ & R8 \\
puszczenie po \SI{0,49}{\second} (z $S_2$) & kliknięcie K -- ekran gaśnie ($S_1$) & R2 \\
puszczenie po \SI{0,51}{\second} (z $S_2$) & kliknięcie Ś -- menu otwarte ($S_{3b}$) & R3 \\
puszczenie po \SI{1,99}{\second} (z $S_2$) & kliknięcie Ś -- menu otwarte ($S_{3b}$) & R3 \\
przytrzymanie \SI{2,01}{\second} (z $S_2$) & kliknięcie D -- telefon wyłączony ($S_0$) & R6 \\
dwa szybkie kliknięcia K (z $S_2$) & ekran gaśnie i zapala się, stan $S_2$ & R2 \\
\bottomrule
\end{tabular}
\end{table}

\subsection{Wyniki}
We wszystkich scenariuszach z tabel~\ref{tab:testy}
i~\ref{tab:graniczne} wyjścia $Q_1$ i $Q_2$ były zgodne z oczekiwanymi
wartościami. Stan $Q_1=1$ przy $Q_2=0$ nigdy nie wystąpił. Przy
trzymaniu przycisku dłużej niż \SI{2}{\second} nie obserwowano już
migotania wyjść, które występowało w pierwszej wersji programu
(punkt~\ref{sec:zdarzenia}).
